% ============================================================================
%  几何内窥与镜像对称
%  Geometric Endoscopy and Mirror Symmetry
%  Edward Frenkel (University of California, Berkeley)
%  Edward Witten (School of Natural Sciences, Institute for Advanced Study)
%  arXiv:0710.5939v3 [math.AG], 2008 年 4 月 5 日, 130 pp.
%
%  中文学术译注本：语义级全译 + 公式逐式核对 + 学术审校批注
%
%  编译：xelatex main.tex（连续两遍，生成目录与交叉引用）
%  中文支持：ctex 的 fandol 字库（跨平台、无需系统字体）
% ============================================================================

% !TEX program = xelatex
\documentclass[11pt,a4paper]{article}
% ============================================================================
%  preamble.tex —— 导言区（宏包、版式、定理环境、记号宏）
%  项目：Frenkel–Witten, Geometric Endoscopy and Mirror Symmetry
%  编译：xelatex main.tex （连续两遍）
% ============================================================================

% !TEX program = xelatex
\usepackage[UTF8,scheme=plain,fontset=fandol]{ctex}

% ---------------------------------------------------------------------------
% 数学与排版
% ---------------------------------------------------------------------------
\usepackage{amsmath}
\usepackage{amssymb}
\usepackage{amsthm}
\usepackage{mathtools}
\usepackage{bm}
% 注意：不载入 amscd，因为宏 \CD 用作 D-模记号；交换图请用 array 环境重建（见 SPEC §4）。
\usepackage[a4paper,top=2.4cm,bottom=2.4cm,left=2.3cm,right=2.3cm]{geometry}
\usepackage{graphicx}
\usepackage{xcolor}
\usepackage{enumitem}
\usepackage{booktabs}
\usepackage{array}
\usepackage{caption}
\usepackage{float}
\usepackage{fancyhdr}
\usepackage{tikz}
\usetikzlibrary{arrows.meta}
\usepackage[hidelinks,bookmarksnumbered=true,bookmarksopen=true]{hyperref}

\hypersetup{
    hypertexnames=false,   % 显式 \tag 编号与自动计数器易产生重名锚点，故关闭
    pdftitle={几何内窥与镜像对称（中文译注本）},
    pdfauthor={Edward Frenkel, Edward Witten（原著）；中文译注},
    pdfsubject={arXiv:0710.5939v3 [math.AG]},
    pdfkeywords={geometric Langlands, endoscopy, mirror symmetry, Hitchin fibration,
                 A-brane, D-module, Hecke eigensheaf, L-packet}
}

% ---------------------------------------------------------------------------
% 版式
% ---------------------------------------------------------------------------
\allowdisplaybreaks
\renewcommand{\contentsname}{目\quad 录}
\renewcommand{\abstractname}{摘\quad 要}
\renewcommand{\refname}{参考文献}
\renewcommand{\figurename}{图}
\renewcommand{\tablename}{表}
\setlength{\emergencystretch}{2em}
\setlength{\parskip}{0.2em}
\linespread{1.06}

\makeatletter
\renewcommand*\l@section{\@dottedtocline{1}{0em}{1.6em}}
\makeatother

% 定理类环境：编号一律照原文手写，避免自动计数与原文不一致
\newenvironment{namedthm}[1]
  {\par\medskip\noindent\textbf{定理 #1.}\itshape\ignorespaces}
  {\par\medskip}
\newenvironment{namedprop}[1]
  {\par\medskip\noindent\textbf{命题 #1.}\itshape\ignorespaces}
  {\par\medskip}
\newenvironment{namedlemma}[1]
  {\par\medskip\noindent\textbf{引理 #1.}\itshape\ignorespaces}
  {\par\medskip}
\newenvironment{namedcor}[1]
  {\par\medskip\noindent\textbf{推论 #1.}\itshape\ignorespaces}
  {\par\medskip}
\newenvironment{namedconj}[1]
  {\par\medskip\noindent\textbf{猜想 #1.}\itshape\ignorespaces}
  {\par\medskip}
\newenvironment{nameddef}[1]
  {\par\medskip\noindent\textbf{定义 #1.}\upshape\ignorespaces}
  {\par\medskip}

% 「本节重点」区块
\newenvironment{keytakeaway}
  {\par\medskip\begin{quote}\small\noindent
   \rule{\linewidth}{0.6pt}\par\nopagebreak\smallskip\noindent
   \textbf{本节重点}\par\nobreak\smallskip\noindent\ignorespaces}
  {\par\smallskip\noindent\rule{\linewidth}{0.6pt}\end{quote}\medskip}

% 「译者说明」区块
\newenvironment{translatornote}
  {\begin{quote}\small\noindent\ignorespaces}
  {\end{quote}\medskip}

% 页眉页脚
\pagestyle{fancy}
\fancyhf{}
\fancyhead[L]{\small 几何内窥与镜像对称}
\fancyhead[R]{\small Frenkel--Witten（中文译注本）}
\fancyfoot[C]{\thepage}
\renewcommand{\headrulewidth}{0.4pt}

% 记号宏（与 GLOSSARY 一致；不要新增）
\newcommand{\Bun}{\mathrm{Bun}}
\newcommand{\CD}{\mathcal{D}}
\newcommand{\CM}{\mathcal{M}}
\newcommand{\CH}{\mathcal{H}}
\newcommand{\CC}{\mathbb{C}}
\newcommand{\ZZ}{\mathbb{Z}}
\newcommand{\RR}{\mathbb{R}}
\newcommand{\PP}{\mathbb{P}}
\newcommand{\Tr}{\operatorname{Tr}}
\newcommand{\tr}{\operatorname{tr}}
\newcommand{\dd}{\mathrm{d}}
\newcommand{\ii}{\mathrm{i}}

% 分片容错：尚未交付时不中断编译
\newcommand{\frag}[1]{%
  \IfFileExists{sections/#1.tex}{\input{sections/#1.tex}}{%
    \par\noindent\textcolor{red}{[缺分片：#1]}\par}}


\begin{document}

\title{\bfseries 几何内窥与镜像对称\\[0.6em]
       \large Geometric Endoscopy and Mirror Symmetry}
\author{Edward Frenkel \quad Edward Witten\\[0.3em]
        \small University of California, Berkeley \quad
        School of Natural Sciences, Institute for Advanced Study}
\date{\small arXiv:0710.5939v3 [math.AG]\\[0.2em]
      \small 2008 年 4 月 5 日}
\maketitle

\begin{abstract}
\noindent
几何 Langlands 对应一直被解释为两个对偶约化群的 Hitchin 纤维化之间的镜像对称。
这一镜像对称又归结为一般 Hitchin 纤维（光滑环面）上的 T-对偶。
本文研究 B-模型一侧的 Hitchin 纤维出现轨形奇点时的情形。
这些奇点对应于自同构群有限的局部系统。
在经典 Langlands 纲领中，这类局部系统称为内窥的。
它们在自守表示理论中，特别是在迹公式的稳定化中，起重要作用。
我们的目标是用 Hitchin 纤维化的镜像对称，揭示这类局部系统在几何理论中所扮演的特殊角色。
对偶 Hitchin 纤维上 A-膜范畴的研究，使我们能揭示几何 Langlands 对应中与内窥相关的一些有趣现象。
随后我们把预言带回到经典的自守函数理论，从而检验并证实它们。
我们所用的几何与 B.-C. Ngô 近期工作中所利用的几何相似，这一事实对理解迹公式可能具有重要意义。
\end{abstract}

\begin{translatornote}
\textbf{译者说明。}
本译本的唯一目标是「无损」：正文逐句对译，公式逐式对照原文页面图像核对，
一切编号（公式 $(x.y)$、文献 $[n]$、章节号、脚注）均保持原样。
原文脚注在译文中保留为脚注，其编号由排版系统重新生成，故与原文页码不同，此为正常现象。

\begin{itemize}[leftmargin=2em,itemsep=1pt,topsep=2pt]
\item \texttt{\% [Tutor Note: \dots]} —— 推导补足、易错预警或物理图像说明；
\item \texttt{\% [原文疑误: \dots]} —— 已判定为原文笔误处，附判断依据与建议；
\item \texttt{\% [存疑: \dots]} —— 无法从图像确证、但疑有问题之处；
\item \texttt{\% [背景补充: \dots]} —— 对中文读者构成实质门槛的术语或背景一句话说明。
\end{itemize}

原文（arXiv:0710.5939v3）全篇 130 页。经位图与矢量图形逐页检测核实：
\textbf{无带框表格、无编号 Figure}；全篇只有一处不编号插图，位于原书 p.18
（奇异 Hitchin 纤维，A-模型 $G=SL_2$ 与 B-模型 ${}^{L}G=SO_3$ 两个图）。
该插图已从原 PDF 按 500 dpi 裁剪为 \texttt{figures/fig-p18-hitchin-fibers.png}，并在 §2 相应位置插入。
文献列表保留英文原貌（书目信息不可翻译）。
\end{translatornote}

\clearpage
\tableofcontents
\clearpage

% ---------------------------------------------------------------------------
% 正文分片（顺序与原文严格一致）
% ---------------------------------------------------------------------------
\frag{s01-intro}
\frag{s02-duality}
\frag{s03a-genus1}
\frag{s03b-genus1}
\frag{s04a-branes}
\frag{s04b-branes}
\frag{s05a-spectral}
\frag{s05b-hecke}
\frag{s06-categories}
\frag{s07-classical}
\frag{s08-functions}
\frag{s09-other}
\frag{s10-gerbes}
\frag{s11-appendix}

% ---------------------------------------------------------------------------
% 参考文献
% ---------------------------------------------------------------------------
\frag{references}

\end{document}
